\subsection{The remaining norm obstruction as a compatible ideal-class family}

\begin{statement}{A.86}{Proved}
{An explicit ideal-class representative of the one-variable obstruction}
Fix a prime $\ell$ and an ordinary pair $(A,B)$ in the scope of
A.77, including its dyadic extension.  Thus every local pair is
$(x_v,1-x_v)$ with $x_v\in\Q_v\setminus\{0,1\}$, and every
completed norm fibre in A.77 is nonempty.  Retain its nested
fields, rational representatives, and compatible generators:
\[
 N_n=\ell^n,\qquad K_n=\Q(\alpha_n),\qquad
 \alpha_n^{N_n}=a_n,\qquad
 a_{n+1}/a_n=b_n^{N_n},\qquad
 \alpha_{n+1}^{\ell}=b_n\alpha_n.
 \tag{A.86.1}
\]
In this statement $b_n$ denotes the rational transition factor,
not a representative of the second completed coordinate.
Write $\widehat{K_n^\times}_{\ell}$ for the inverse limit of
$K_n^\times/(K_n^\times)^{\ell^j}$.

There are explicit norm-compatible completed divisors
$\mathfrak d_n$ with norm equal to the divisor of $B$.
For odd $\ell$ these are finite-prime divisors; for $\ell=2$
they also retain the signs at all real embeddings.
Let $\mathcal D_n$ be the corresponding completed divisor group,
and let $\mathcal C_n$ be the $\ell$-primary ordinary ideal class
group for odd $\ell$, and the two-primary narrow ideal class
group for $\ell=2$.  Define
\[
\begin{aligned}
 W_n&=\ker\bigl(\mathrm N_{n,0}:\mathcal D_n\longrightarrow
                                      \mathcal D_0\bigr),\\
 C_n^0&=\operatorname{im}(W_n\longrightarrow\mathcal C_n),\\
 Z_n&=\ker\bigl(\mathrm N_{n,0}:\widehat{K_n^\times}_{\ell}
                         \longrightarrow\widehat{\Q^\times}_{\ell}\bigr).
\end{aligned}
 \tag{A.86.2}
\]
All transition maps are field norms or ideal norms.  The family
$c_n=[\mathfrak d_n]$ is an element of $\varprojlim_n C_n^0$.
There is a natural exact sequence
\[
 \varprojlim_n W_n\longrightarrow\varprojlim_n C_n^0
 \xrightarrow{\partial_{\rm cl}}\varprojlim_n{}^1 Z_n
 \longrightarrow\varprojlim_n{}^1 W_n\longrightarrow0.
 \tag{A.86.3}
\]
For the sign convention of A.81, its injection
$\iota:\varprojlim{}^1 Z_n\hookrightarrow
H^2_{\rm cts}(G_\Q,D_A)$ satisfies
\[
             \iota\bigl(\partial_{\rm cl}((c_n)_n)\bigr)
                             =o_{\rm pol}(A,B).
 \tag{A.86.4}
\]
Consequently the exact full one-variable lifting condition is
\[
 (c_n)_n\in
 \operatorname{im}\left(\varprojlim_n W_n
                              \longrightarrow\varprojlim_n C_n^0\right).
 \tag{A.86.5}
\]
In particular the actual obstruction has zero image in
$\varprojlim{}^1 W_n$.  Formula (A.86.5) retains a compatibility
problem in finite ideal class groups and their completed divisor
representatives; this statement does not assert that the indicated
family belongs to that image.
\end{statement}
\begin{proof}
We first give the exact completed divisor sequences, then construct
the particular compatible divisors from the ordinary local points.

\emph{Divisors, signs, and compact units.}
For a number field $K$, put
\[
 \widehat{\operatorname{Div}}_\ell(K)
 =\varprojlim_j\bigoplus_{\mathfrak p}
                                  \Z/\ell^j\Z\,[\mathfrak p]
 =\left\{(m_{\mathfrak p})_{\mathfrak p}\in
       \prod_{\mathfrak p}\Z_\ell:
       m_{\mathfrak p}\longrightarrow0\text{ $\ell$-adically}\right\}.
 \tag{A.86.6}
\]
The limit in the last expression means that for every $j$ only
finitely many entries are nonzero modulo $\ell^j$.
For odd $\ell$ take $\mathcal D(K)=
\widehat{\operatorname{Div}}_\ell(K)$.
For $\ell=2$ take
$\mathcal D(K)=\widehat{\operatorname{Div}}_2(K)
\oplus(\Z/2\Z)^{\operatorname{Hom}(K,\R)}$,
and map an element of $K^\times$ to its finite valuations and
its real signs, with a negative sign represented by one.
The ideal norm is
$[\mathfrak P]\mapsto f(\mathfrak P/\mathfrak p)[\mathfrak p]$.
On signs it is the sum over real embeddings above a fixed real
embedding; pairs of complex embeddings contribute a positive norm.
These formulas define continuous norm maps on $\mathcal D(K)$.

For odd $\ell$, let $\mathcal U(K)$ be the $\ell$-adic completion
of $\mathcal O_K^\times$.  For $\ell=2$, let it be the two-adic
completion of the totally positive units.  There is an exact
sequence
\[
 0\longrightarrow\mathcal U(K)\longrightarrow
 \widehat{K^\times}_\ell\xrightarrow{\operatorname{div}^{\rm sgn}}
 \mathcal D(K)\longrightarrow\mathcal C(K)\longrightarrow0.
 \tag{A.86.7}
\]
Here $\mathcal C(K)$ is the indicated ordinary or narrow primary
class group.  Before completion the dyadic class map sends
$(D,\varepsilon)$ to $[D-\operatorname{div}(\gamma)]^+$,
where $\gamma\in K^\times$ has sign vector $\varepsilon$.
Weak approximation supplies $\gamma$; two choices have totally
positive quotient, so the class is well-defined.  Its kernel is
exactly the principal signed divisors.  Thus the quotient of
the signed divisor group by principal signed divisors is then
the quotient of finite divisors by principal divisors generated
by totally positive elements, which is exactly the narrow class
group.  Its finiteness, and the finite generation of the unit
group, give (A.86.7) after completion: a subgroup of finite index
in the divisor-and-sign group has the induced $\ell$-adic topology,
and the principal-divisor group is a free abelian group with at
most the finite sign torsion.  Equivalently, taking the finite
power-quotient sequences gives (A.86.7), since the inverse limits
of the finite class-group torsion kernels have zero Tate module.
For odd $\ell$ the same argument uses the ordinary class group.

The groups $\mathcal U(K)$ are compact finitely generated
$\Z_\ell$-modules.  In the signed convention
$\mathcal U(\Q)=0$ and $\mathcal C(\Q)=0$ for every $\ell$.
Thus $\operatorname{div}^{\rm sgn}$ identifies
$\widehat{\Q^\times}_\ell$ with $\mathcal D_0$ exactly.
In particular a principal divisor in $W_n$ has a preimage with
completed field norm one, rather than merely a norm with zero
finite valuations.  Taking kernels in (A.86.7) therefore gives
\[
 0\longrightarrow\mathcal U(K_n)\longrightarrow Z_n
       \longrightarrow W_n\longrightarrow C_n^0\longrightarrow0.
 \tag{A.86.8}
\]
Let $P_n^0$ be the image of $Z_n$ in $W_n$.
Every countable inverse system of compact abelian groups has
zero algebraic $\varprojlim{}^1$: the finite sets of equations
$u_n-f_nu_{n+1}=v_n$ have solutions by backward substitution,
and compact finite-intersection gives a simultaneous solution.
Apply this to the completed units in (A.86.8).  The two short
exact sequences obtained from (A.86.8) give
$\varprojlim{}^1Z_n\simeq\varprojlim{}^1P_n^0$ and (A.86.3).
Here a countable inverse system has no $\varprojlim{}^2$, and
$\varprojlim{}^1 C_n^0=0$ because the groups $C_n^0$ are finite.

\emph{Compatible local norm elements.}
For each place $v$, ordinary equality with $A$ gives
$x_v/a_n\in(\Q_v^\times)^{N_n}$.
Choose roots $z_{n,v}\in\Q_v^\times$ such that
\[
 z_{n,v}^{N_n}=x_v/a_n,\qquad
 z_{n,v}=b_nz_{n+1,v}^{\ell},\qquad
 \gamma_{n,v}=1-z_{n,v}\alpha_n
                         \in(K_n\otimes_\Q\Q_v)^\times.
 \tag{A.86.9}
\]
Compatible roots exist without an extra divisibility assumption.
At each level the root set is finite and nonempty.  The displayed
transition sends a root at level $n+1$ to a root at level $n$;
every finite string is obtained from a root at its highest level.
Compactness of the product of these finite sets gives an infinite
compatible string.  The same argument applies at the real place.
As $x_v\ne1$, every $\gamma_{n,v}$ is invertible.  The norm
polynomial for $T^\ell-b_n\alpha_n$ gives the literal equalities
\[
 \mathrm N_{K_{n+1}\otimes\Q_v\,/\,K_n\otimes\Q_v}
                               (\gamma_{n+1,v})=\gamma_{n,v},
 \qquad
 \mathrm N_{K_n\otimes\Q_v\,/\,\Q_v}(\gamma_{n,v})=1-x_v.
 \tag{A.86.10}
\]
Thus their finite valuations, and their real signs when $\ell=2$,
are norm-compatible.  It remains to prove that these valuations
belong to the completed global divisor group (A.86.6).

Fix $n$ and exclude the finite set consisting of $\ell$ and the
prime divisors of the numerator and denominator of $a_n$.
At every remaining prime $q$, the algebra $K_n\otimes\Q_q$ is
unramified, since $T^{N_n}-a_n$ has unit discriminant.
The element $\alpha_n$ is a unit in each factor.
Write $m=v_q(x_q)$ and $t=v_q(1-x_q)$; both are ordinary integers.
The ordinary unit equation gives exactly the following four cases.

If $m>0$ and $t=0$, then $v_q(z_{n,q})=m/N_n>0$, so every
factor of $\gamma_{n,q}$ is a unit and contributes zero.
If $m=t=0$, both $z_{n,q}$ and $\alpha_n$ are units, and all
factors of $\gamma_{n,q}$ are integral.  Their norm is a unit,
so every factor is a unit.
If $m=0$ and $t>0$, reduction of $\gamma_{n,q}$ can vanish
only at the simple linear factor with
$\overline\alpha_n=\overline z_{n,q}^{-1}$.
There is precisely one such factor: the polynomial is separable,
and $x_q\equiv1\pmod q$ makes that value a root.
Its residue degree is one; all other factors are units.  The
norm equality then gives valuation exactly $t$ at that branch.
Finally, if $m=t<0$, then $v_q(z_{n,q})=m/N_n<0$, and in every
unramified factor the valuation of $1-z_{n,q}\alpha_n$ is
exactly $m/N_n$.

The valuation vectors of the completed classes $A$ and $B$
tend to zero $\ell$-adically.  The preceding four formulas show
that every entry just constructed is either zero, $t$, or
$m/N_n$.  Division by the one fixed integer $N_n$ preserves
the convergence to zero.  Each rational prime has at most
$N_n$ primes above it.  Thus the entire family over the finite
primes of $K_n$ belongs to (A.86.6); the excluded finite set
does not change this conclusion.  Add the finitely many real
signs of $\gamma_{n,\infty}$ when $\ell=2$, and call the result
$\mathfrak d_n$.  Equation (A.86.10) proves compatibility and
$\mathrm N_{n,0}\mathfrak d_n=\operatorname{div}^{\rm sgn}(B)$.

\emph{The class family and its sign in the obstruction group.}
By A.77 choose any $\beta_n\in\widehat{K_n^\times}_\ell$
with $\mathrm N_{n,0}\beta_n=B$.  No compatibility is assumed.
Put
\[
 w_n=\mathfrak d_n-\operatorname{div}^{\rm sgn}(\beta_n)
                \in W_n,\qquad
 d_n'=\beta_n/\mathrm N_{n+1,n}(\beta_{n+1})\in Z_n.
 \tag{A.86.11}
\]
The class of $w_n$ is $c_n=[\mathfrak d_n]$, so $c_n\in C_n^0$.
The compatibility of $\mathfrak d_n$ makes $(c_n)_n$ a
compatible family.  For the difference convention
$\Delta(w)_n=w_n-\mathrm N_{n+1,n}w_{n+1}$, (A.86.11) gives
\[
                  \Delta(w)_n
                         =-\operatorname{div}^{\rm sgn}(d_n').
 \tag{A.86.12}
\]
Consequently $\partial_{\rm cl}((c_n)_n)$ corresponds to
$-[d_n']$ under the isomorphism
$\varprojlim{}^1 Z_n\simeq\varprojlim{}^1P_n^0$.
A.81 gives $\iota([d_n'])=-o_{\rm pol}(A,B)$, proving
(A.86.4).  The injection $\iota$ and exactness of (A.86.3)
now prove (A.86.5), with its actual source and target maps.

One can also recover the lifting implication directly.
A compatible family in $W_n$ representing $(c_n)_n$ makes
$\mathfrak d_n-w_n$ a compatible family of completed principal
signed divisors of norm $B$.  Choose principal generators.
Their successive norm discrepancies are completed units;
$\varprojlim{}^1\mathcal U(K_n)=0$ corrects those generators
to a norm-compatible family.  A.77 then gives the full actual
one-variable section.  Conversely any norm-compatible family
of norm $B$ subtracts from $(\mathfrak d_n)_n$ to give exactly
such a family in $\varprojlim W_n$.
\end{proof}

The construction localizes the remaining compatibility question
to the specific finite ideal-class family $([\mathfrak d_n])_n$.
A full decomposition group at one place proves the finite-level
Hasse norm statements used in A.77.  No surjectivity of
$\varprojlim W_n\to\varprojlim C_n^0$ has been deduced from
that fact here.

\noindent The exact certificate is
\path{code/norm_ideal_class_certificate.py}, with retained output
\path{code/norm_ideal_class_certificate_output.txt}.
It checks the literal norm polynomials, the unique degree-one
residue branch, the fixed-denominator convergence bounds, and
the connecting sign.  It makes no ideal-class vanishing claim.
